\section{Introduction}
\label{sec:intro}

% Intended content (~2 pages):
%
%  1.1 Fisher's paired efficiency and its generalisation.
%      One paragraph recapping Fisher's 1/(1 - rho) and the PEF generalisation
%      eta = (1 + kappa)/(1 + kappa - 2*sqrt(kappa)*rho).  Cite the empirical
%      companion paper for the cross-domain validation; this paper studies
%      the structure of the formula itself.
%
%  1.2 Why the formula has more structure than it appears to.
%      Motivate the half-strip reparameterisation tau = (1/2) log kappa and
%      the canonical form eta = cosh tau / (cosh tau - rho).  Explain that
%      everything in Sections 2--6 is a consequence of this single identity.
%
%  1.3 Contributions.
%      List six contributions:
%        (i)   canonical form on the half-strip and the kappa <-> 1/kappa
%              involution;
%        (ii)  Moebius linearisation 1/eta = 1 - rho sech tau and the
%              Chebyshev/Poisson series;
%        (iii) Riemann-sphere realisation eta = 1/(1 - cos sigma) and the
%              identification of the Fisher regime (kappa = 1) with the
%              sphere's equator;
%        (iv)  partition-function reading log eta = -log(1 - rho sech tau)
%              and the latent geometric exponential family;
%        (v)   Fisher-Rao geometry, the variance-stabilising coordinate psi,
%              and the regime change at rho = 0;
%        (vi)  numerical validation of (i)--(v) on the six datasets of the
%              empirical companion.
%
%  1.4 Relation to the empirical companion paper.
%      Forward reference: empirical companion uses Sections 2--6 results as
%      cited infrastructure; this paper uses Section 7's numerical results to
%      ground the geometry.  Each paper is independently readable.
%
%  1.5 Outline.
%      Pointer to Sections 2--8.

% TODO: drafting.
